% !TEX program = pdflatex
% Plantilla base · pdfLaTeX · català
\documentclass[a4paper,11pt]{article}

% --- Codificació i llengua (pdfLaTeX) -------------------------------
\usepackage[T1]{fontenc}
\usepackage[utf8]{inputenc}
\usepackage[catalan]{babel}
\usepackage{lmodern}
\usepackage{microtype}

% --- Pàgina ---------------------------------------------------------
\usepackage[a4paper,margin=2.5cm]{geometry}

% --- Text, citacions i números --------------------------------------
\usepackage{csquotes}     % \enquote{...}
% csquotes no porta estil català: el definim (« » exteriors, “ ” interiors)
\DeclareQuoteStyle{catalan}
  {\guillemotleft}{\guillemotright}
  {\textquotedblleft}{\textquotedblright}
\usepackage{siunitx}      % números i unitats amb coma decimal
\sisetup{output-decimal-marker={,}, group-separator={\,}}

% --- Matemàtiques ---------------------------------------------------
\usepackage{amsmath,amssymb,amsthm}
\DeclareMathOperator{\tg}{tg}       % tangent (notació catalana)
\DeclareMathOperator{\arctg}{arctg}

\theoremstyle{plain}
\newtheorem{teorema}{Teorema}
\newtheorem{lema}[teorema]{Lema}
\newtheorem{proposicio}[teorema]{Proposició}
\newtheorem{corollari}[teorema]{Corol·lari}
\theoremstyle{definition}
\newtheorem{definicio}[teorema]{Definició}
\newtheorem{exemple}[teorema]{Exemple}
\theoremstyle{remark}
\newtheorem*{observacio}{Observació}

% --- Figures, taules i enllaços -------------------------------------
\usepackage{graphicx}
\usepackage{booktabs}
\usepackage[hidelinks]{hyperref}    % carregar-lo gairebé al final

% --- Bibliografia (descomenteu si cal; requereix executar biber) ----
% \usepackage[backend=biber,style=numeric]{biblatex}
% \addbibresource{referencies.bib}

\title{Títol del document}
\author{Nom Cognom}
\date{\today}

\begin{document}
\maketitle

\begin{abstract}
Resum breu del contingut del document.
\end{abstract}

\tableofcontents

\section{Introducció}
\label{sec:introduccio}

Text d'exemple amb una citació: \enquote{això és una cita}. A la
secció~\ref{sec:resultats} es presenten els resultats; el valor és
\num{3,14} i la massa, \SI{2,5}{\kilo\gram}.

\section{Resultats}
\label{sec:resultats}

\begin{teorema}\label{thm:principal}
Per a tot $x\in\mathbb{R}$ es compleix $\sin^2 x + \cos^2 x = 1$.
\end{teorema}
\begin{proof}
Conseqüència directa del teorema de Pitàgores.
\end{proof}

\begin{table}[ht]
  \centering
  \caption{Exemple de taula.}
  \label{tab:exemple}
  \begin{tabular}{lr}
    \toprule
    Concepte & Valor \\
    \midrule
    Primer   & \num{1,5} \\
    Segon    & \num{2,0} \\
    \bottomrule
  \end{tabular}
\end{table}

La taula~\ref{tab:exemple} resumeix les dades.

% \printbibliography

\end{document}
